\begin{lemma}
\label{lemma-points}
Let $S$ be a scheme. Let $(U, R, s, t, c)$ be a groupoid scheme over $S$.
Assume $U = \Spec(A)$ and $R = \Spec(B)$ are affine and
$s, t : R \to U$ finite locally free. Let $C \subset A$ be as in
(\ref{equation-invariants}). Then $U \to M = \Spec(C)$ has
the following properties:
\begin{enumerate}
\item the map on points $|U| \to |M|$ is surjective and
$u_0, u_1 \in |U|$ map to the same point if and only if
there exists a $r \in |R|$ with $t(r) = u_0$ and $s(r) = u_1$, in
a formula
$$
|M| = |U|/|R|
$$
\item for any algebraically closed field $k$ we have
$$
M(k) = U(k)/R(k)
$$
\end{enumerate}
\end{lemma}

\begin{proof}
Since $C \to A$ is integral (Lemma \ref{lemma-integral-over-invariants})
and injective we see that $\Spec(A) \to \Spec(C)$ is surjective, see
Algebra, Lemma \ref{algebra-lemma-integral-overring-surjective}.
Thus $|U| \to |M|$ is surjective.

\medskip\noindent
Let $k$ be an algebraically closed field and let $C \to k$ be a ring map.
Since surjective morphisms are preserved under base change
(Morphisms, Lemma \ref{morphisms-lemma-base-change-surjective})
we see that $A \otimes_C k$ is not zero. Now $k \subset A \otimes_C k$ is a
nonzero integral extension. Hence any residue field of $A \otimes_C k$
is an algebraic extension of $k$, hence equal to $k$. Thus we see that
$U(k) \to M(k)$ is surjective.

\medskip\noindent
Let $a_0, a_1 : A \to k$ be two ring maps. If there exists a ring map
$b : B \to k$ such that $a_0 = b \circ t^\sharp$ and $a_1 = b \circ s^\sharp$
then we see that $a_0|_C = a_1|_C$ by definition. Thus the map
$U(k) \to M(k)$ equalizes the two maps $R(k) \to U(k)$.
Conversely, suppose that $a_0|_C = a_1|_C$. Let us name this algebra
map $c : C \to k$. Consider the diagram
$$
\xymatrix{
& &
B \ar@{-->}[lld] \\
k & &
A
\ar@<0.5ex>[ll]^{a_0}
\ar@<-0.5ex>[ll]_{a_1}
\ar@<1ex>[u]
\ar@<-1ex>[u] \\
& &
C \ar[u] \ar[llu]^c
}
$$
If we can construct a dotted arrow making the diagram commute, then
the proof of part (2) of the lemma is complete. Since $s : A \to B$ is finite
there exist finitely many ring maps
$b_1, \ldots, b_n : B \to k$ such that $b_i \circ s^\sharp = a_1$.
If the dotted arrow does not exist, then we see that none of the
$a'_i = b_i \circ t^\sharp$, $i = 1, \ldots, n$ is equal to $a_0$.
Hence the maximal ideals
$$
\mathfrak m'_i = \Ker(a_i' \otimes 1 : A \otimes_C k \to k)
$$
of $A \otimes_C k$ are distinct from
$\mathfrak m = \Ker(a_0 \otimes 1 : A \otimes_C k \to k)$.
By Algebra, Lemma \ref{algebra-lemma-silly} we would get an element
$f \in A \otimes_C k$ with $f \in \mathfrak m$, but
$f \not \in \mathfrak m_i'$ for $i = 1, \ldots, n$.
Consider the norm
$$
g = \text{Norm}_{s^\sharp \otimes 1}(t^\sharp \otimes 1(f))
\in
A \otimes_C k
$$
By Lemma \ref{lemma-determinant-trick} this lies in the invariants
$C^1 \subset A \otimes_C k$ of the base change
groupoid (base change via the map $c : C \to k$). On the one hand,
$a_1(g) \in k^*$ since
the value of $t^\sharp(f)$ at all the points (which correspond to
$b_1, \ldots, b_n$) lying over $a_1$ is
invertible (insert future reference on property determinant here).
On the other hand, since $f \in \mathfrak m$, we see that
$f$ is not a unit, hence $t^\sharp(f)$ is not a unit
(as $t^\sharp \otimes 1$ is faithfully flat),
hence its norm is not a unit (insert future reference
on property determinant here). We conclude that $C^1$ contains
an element which is not nilpotent
and not a unit. We will now show that this leads to a contradiction.
Namely, apply Lemma \ref{lemma-invariants-base-change}
to the map $c : C \to C' = k$, then
we see that the map of $k$ into the invariants $C^1$ is injective
and moreover, that for any element $x \in C^1$ there exists an integer
$n > 0$ such that $x^n \in k$. Hence every element of $C^1$ is
either a unit or nilpotent.

\medskip\noindent
We still have to finish the proof of (1). We already know that
$|U| \to |M|$ is surjective. It is clear that $|U| \to |M|$ is
$|R|$-invariant. Finally, suppose $u_0, u_1 \in U$ maps to the same
point $m \in M$. Then the induced field extensions $\kappa(u_0)/\kappa(m)$
and $\kappa(u_1)/\kappa(m)$ are algebraic (as $A$ is integral over $C$
as used above). Hence if $k$ is an algebraic closure of $\kappa(m)$
then we can find $\kappa(m)$-embeddings $\overline{u}_0 : \kappa(u_0) \to k$
and $\overline{u}_1 : \kappa(u_1) \to k$. These determine $k$-valued
points $\overline{u}_0, \overline{u}_1 \in U(k)$ mapping to the same
point of $M(k)$. By part (2) we see that there exists a point
$\overline{r} \in R(k)$ with $s(\overline{r}) = \overline{u}_0$ and
$t(\overline{r}) = \overline{u}_1$. The image $r \in R$ of $\overline{r}$
is a point with $s(r) = u_0$ and $t(r) = u_1$ as desired.
\end{proof}
